\section{Introduction}

Procurement set-asides for small firms are among the most widely deployed government-procurement instruments worldwide \citep{bosio2022}. They were designed on a textbook intuition: bringing more bidders to a protected segment should reduce procurement prices. The intuition has not survived contact with the data. A decade of empirical evidence places the price cost of bidder-exclusion rules in the ten-to-fifteen percent range \citep{marion2007, krasnokutskaya2011, athey2013, nakabayashi2013, hyytinen2018}. Two non-exclusive mechanisms operate behind that markup: the bidder pool shrinks, and the auction clears at a different point in the surviving cost distribution---an order-statistic shift the structural framework identifies separately from any active behavioral adjustment by the surviving bidders. Price regressions identify the total but cannot disentangle the two, and the welfare ranking of bidder exclusion against alternative preference instruments turns on the decomposition---a ranking that, as this paper shows, is conditional on the thickness of the protected pool and the heterogeneity of the good rather than universal. On São Paulo's medical-supplies procurement alone, the structural exercise prices the rule at R\$\,\welfAnnualBRLLo--\welfAnnualBRLHi~million (US\$\,\welfAnnualUSDLo--\welfAnnualUSDHi~million) per year and decomposes the loss into the two channels.

A \policyCutoffMonth{} administrative reversal in S\~ao Paulo state, Brazil, supplies the natural experiment. The state extended an existing SME-only procurement rule to medical and hospital supplies---the last major exception---on the Bolsa Eletr\^onica de Compras (BEC), the country's largest centralized state-government procurement platform. The reversal was a legal opinion internal to the state prosecutor's office, defended on isonomy grounds and uncorrelated with prices, competition, or supplier complaints in the affected product group; the timing was bureaucratic. A difference-in-differences against \controlGroupCount{} never-treated product groups (\ref{app:did}) finds that absolute winning prices rose by nearly \didLowerPct~percent and SME participation roughly doubled. The DiD coefficient fixes the average effect on observed log prices; the structural exercise of this paper recovers which mechanism produced the markup and what the rule costs in welfare terms, anchored to the same magnitude.

I specify an asymmetric independent-private-values reading of the platform's Preg\~ao English-reverse auction format: bidders draw costs from type-specific distributions $F_c^{\text{SME}}$ and $F_c^{\neg}$ and compete under a descending clock. Drop-out prices of losers therefore point-identify $F_c^k$ under \citet{haile2003}, without inverting a bid function or imposing parametric structure. An auction-level common scale is shrunk out following \citet{krasnokutskaya2011}; winners are treated as left-censored at the second-lowest exit and recovered by a Turnbull NPMLE robustness. Entry is handled in the \citet{atheyseira2011} spirit: pre- and post-period bidder counts as equilibrium arrival rates, with the main specification allowing the post-policy SME cost distribution to differ from its pre-policy counterpart as an equilibrium-selection object. The exercise is conducted under observed equilibrium entry; it is not a fully solved nested entry-bidding equilibrium, and the decomposition is a counterfactual-simulation accounting under the modeling assumptions.\footnote{Three magnitudes appear in this paper and should not be conflated. The \emph{DiD coefficient} (\didLowerPct{}--\didUpperPct{} percent) is the average effect on observed log winning prices net of item and month fixed effects. The \emph{structural counterfactual} ($+\bneEffectTotalNp$ in non-pharma, $+\bneEffectTotalPh$ in pharma) is the simulated mean of the second-order statistic, scaled by the reference price; in $p_{S_1}$ units it is two to three times the DiD coefficient because the two objects average over different primitives. The \emph{welfare cost} (\welfLossPctLthirtyNp{}/\welfLossPctLthirtyPh{}~percent of $p_{S_1}$ at $\lambda = \lambdaMain$) adds the MCPF distortion to the allocative wedge.}

Three findings organize the paper. \emph{(i) The price effect operates predominantly within the auction---sheltered bidding.} Letting $p_{S_1}, p_{S_2}, p_{S_3}$ denote the simulated winning prices under the open regime, the SME-only counterfactual at the pre-policy SME pool, and the SME-only counterfactual at the observed post-policy SME pool, the within-auction piece $p_{S_2} - p_{S_1}$ is \bneShareIntensNp~percent of the total effect in non-pharma and \bneShareIntensPh~percent in pharma (Table~\ref{tab:v3_bne_decomp}). The share stays in the \withinShareLow--\withinShareHigh~percent range across estimators (Turnbull NPMLE: \turnbullShareNp/\turnbullSharePh; strict invariance: \strictShareNp/\strictSharePh). \emph{(ii) Endogenous SME entry is a partial offset, not the source of the markup.} Holding the SME pool at its pre-policy size while imposing the SME-only rule, the simulated price effect rises to \latentShockTimesNp/\latentShockTimesPh{}~percent of the realized one, so without endogenous entry the rule would cost \latentShockRatioNp/\latentShockRatioPh{}~percent more. \emph{(iii) A \optPrefThreshold{} price preference welfare-dominates the set-aside in thick standardized markets at near-zero fiscal cost; pharma bifurcates; the empate ficto sits even closer to the open regime.} In thick standardized markets the simulated price shifts by $\optPrefShiftNp$ under the preference, against $+\bneEffectTotalNp$ under the set-aside, with dominance preserved under strict invariance. In thin heterogeneous pharma the ranking flips under strict invariance, with the implicit welfare weight falling to \welfWeightSiPh. A structural quantification of the V4 \emph{empate ficto} right-to-match operative in current Brazilian law (Table~\ref{tab:v4_empate}) places it firmly below V3 on both margins: a gain of \vFourSmeWinGainNp/\vFourSmeWinGainPh{}~pp in non-pharma/pharma at a welfare cost of \vFourLossPctNp/\vFourLossPctPh{}~percent of $p_{S_1}$. The bifurcation between V0 and the price-preference instruments is the diagnostic.

The paper makes three contributions. \emph{First}, a within-auction structural decomposition of an SME procurement preference identified from a single legal trigger, on a combination of conditions rare in the literature: total-exclusion trigger, Preg\~ao drop-out point-identification, auction-level UH correction \citep{krasnokutskaya2011}, observed-equilibrium-entry treatment \`a la \citet{atheyseira2011}, and a \citet{saezstantcheva2016} welfare-weight identity. \citet{marion2007} estimates the price effect of federal SME set-asides via reduced-form variation; \citet{krasnokutskaya2011} decomposes California highway bid preferences via cross-equation constraints; \citet{athey2013} prices the efficiency-vs-redistribution trade-off between set-asides and subsidies in U.S.\ timber. This paper extends the second approach to a total-exclusion rule on a clean legal trigger, identifying the market-structure conditions under which the standard preference-over-exclusion ranking does not transport. \emph{Second}, endogenous SME entry is part of the welfare object, not a side issue: a fixed-pool counterfactual misstates the realized policy cost by roughly half. \emph{Third}, the welfare comparison is a market-design statement, not a universal ranking---the \optPrefThreshold{} preference welfare-dominates in thick standardized markets, while pharma's bifurcation identifies precisely the market conditions under which the modeling treatment of the protected pool is first-order. The paper sits in dialogue with three lines of empirical work: emerging-market procurement design \citep{decarolis2024, coviello2026, ferraz2016, szerman2023}, the political economy of public contracting \citep{colonnelli2022}, and the public-finance evaluation of government programs in developing economies \citep{best2023, bandiera2021, olken2007}.

A pair of furosemide 40~mg purchases by the same São Paulo public buyer eight months apart and bracketing the cutoff fixes the magnitudes. In February 2018 the unit cleared roughly \furosBelowRefPct~percent below the buyer's reference price; in October 2018 just above. The difference is one realization of the average direction the design estimates across the \nObsAuctions{} firm-auction observations on each side of the \policyCutoffMonth{} cutoff. Table~\ref{tab:preview} summarizes the headline numbers across the four reporting layers used throughout the paper; the rows are not alternative estimates of the same object but distinct objects identified at distinct levels of the framework, related by the mechanical bridge in Section~\ref{sec:bridge} and the welfare-weight identity of Section~\ref{sec:discussion}.

\begin{table}[!htbp]
\centering
\footnotesize
\setlength{\tabcolsep}{4pt}
\begin{threeparttable}
\caption{Headline magnitudes.}
\label{tab:preview}
\begin{tabular}{@{}p{0.65\linewidth} c c@{}}
\toprule
 & \textbf{Non-pharma} & \textbf{Pharma} \\
\midrule
DiD on $\log p^{\mathrm{final}}$ (absolute, App.~\ref{app:did}) & \multicolumn{2}{c}{$+\didLowerDec$ to $+\didUpperDec$} \\
DiD on $p^{\mathrm{final}}/p^{\mathrm{ref}}$ (App.~\ref{app:did}) & \multicolumn{2}{c}{$+\didRefDec$} \\
Simulated $p_{S_3} - p_{S_1}$, units of $p^{\mathrm{ref}}$ & $+\bneEffectTotalNp$ & $+\bneEffectTotalPh$ \\
\addlinespace
Within-auction share, main specification & \bneShareIntensNp\% & \bneShareIntensPh\% \\
Within-auction share, Turnbull NPMLE & \turnbullShareNp\% & \turnbullSharePh\% \\
Within-auction share, strict primitive invariance & \strictShareNp.0\% & \strictSharePh.0\% \\
\addlinespace
Welfare cost (\% of $p_{S_1}$, $\lambda = \lambdaMain$) & \welfLossPctLthirtyNp\% & \welfLossPctLthirtyPh\% \\
$w^{\text{SME}}_\star$ to prefer set-aside, main spec. & \welfWeightMainNp & \welfWeightMainPh \\
$w^{\text{SME}}_\star$, strict invariance & \welfWeightSiNp & \welfWeightSiPh \\
\bottomrule
\end{tabular}
\begin{tablenotes}[flushleft]\footnotesize
\item DiD magnitudes from \ref{app:did} (against \controlGroupCount{} never-treated product groups, \structuralWindowM-month symmetric window around \policyCutoffMonth). Structural magnitudes from the asymmetric IPV counterfactual under observed equilibrium entry; sources by row: Sec.~\ref{sec:results} (decomposition), Sec.~\ref{sec:robustness} (Turnbull NPMLE; strict invariance), Sec.~\ref{sec:welfare} (welfare cost), Sec.~\ref{sec:discussion} (welfare weights). The pharmaceutical V3-vs-V0 ranking reverses under strict primitive invariance ($w^{\text{SME}}_\star = \welfWeightSiPh < 1$); the non-pharmaceutical ranking is preserved under both specifications.
\end{tablenotes}
\end{threeparttable}
\end{table}
